\documentclass[11pt,a4paper]{article}
% Preambulo comun de los reportes de Tecnologias Celulares.
% Cada main.tex hace: \documentclass ... % Preambulo comun de los reportes de Tecnologias Celulares.
% Cada main.tex hace: \documentclass ... % Preambulo comun de los reportes de Tecnologias Celulares.
% Cada main.tex hace: \documentclass ... \input{../comun/preambulo.tex}
\usepackage[utf8]{inputenc}
\usepackage[T1]{fontenc}
\usepackage[spanish,mexico]{babel}
\usepackage{graphicx}
\usepackage[margin=2.5cm]{geometry}
\usepackage{amsmath}
\usepackage{amssymb}
\usepackage{booktabs}
\usepackage{tabularx}
\usepackage{array}
\usepackage{float}
\usepackage{setspace}
\usepackage{parskip}
\usepackage{fancyhdr}
\usepackage{textcomp}
\usepackage{xcolor}
\usepackage{caption}
\usepackage{fvextra}
\usepackage[hidelinks]{hyperref}

\onehalfspacing
\setlength{\parskip}{6pt}
\sloppy
\setlength{\emergencystretch}{3em}
\setlength{\headheight}{40pt}
\captionsetup{font=small,labelfont=bf}
\setlength{\tabcolsep}{6pt}

\pagestyle{fancy}
\fancyhf{}
\fancyhead[L]{\includegraphics[width=1cm]{/home/yovick/repos/UAZ-IEI/Assets/UIE.png}}
\fancyhead[R]{\small Tecnologías Celulares}
\fancyfoot[C]{\thepage}
\renewcommand{\headrulewidth}{0.4pt}

\DefineVerbatimEnvironment{codigo}{Verbatim}{%
  breaklines=true,breakanywhere=true,fontsize=\small,frame=single,
  framesep=3pt,rulecolor=\color{gray!50}}

% Figura centrada con pie: \figura[ancho]{archivo}{leyenda}
\newcommand{\figura}[3][0.85]{%
  \begin{center}
    \includegraphics[width=#1\linewidth]{#2}\\[2pt]
    {\small #3}
  \end{center}}

% Portada: \portada{numero}{titulo de la practica}
\newcommand{\portada}[2]{%
\begin{titlepage}
  \thispagestyle{empty}
  \begin{center}
    \vspace*{1cm}
    \includegraphics[width=4.5cm]{/home/yovick/repos/UAZ-IEI/Assets/uaz.png}
    \vspace{0.8cm}

    {\large Universidad Autónoma de Zacatecas}\\[4pt]
    {\large Unidad Académica de Ingeniería Eléctrica}\\[4pt]
    {\normalsize Ingeniería Electrónica Industrial con especialidad en Telecomunicaciones}

    \vspace{1.8cm}
    {\LARGE\bfseries Práctica #1}\\[10pt]
    {\Large\bfseries #2}
  \end{center}

  \vspace{2.2cm}
  \begin{center}
    \begin{tabular}{rl}
      \textbf{Materia:} & Tecnologías Celulares\\
      \textbf{Docente:} & Dr. Víktor Iván Rodríguez Abdalá\\
      \textbf{Alumno:}  & Arnold Yovick Ríos Zaldívar\\
      \textbf{Fecha:}   & Noviembre de 2026\\
    \end{tabular}
  \end{center}
  \vfill
\end{titlepage}}

\usepackage[utf8]{inputenc}
\usepackage[T1]{fontenc}
\usepackage[spanish,mexico]{babel}
\usepackage{graphicx}
\usepackage[margin=2.5cm]{geometry}
\usepackage{amsmath}
\usepackage{amssymb}
\usepackage{booktabs}
\usepackage{tabularx}
\usepackage{array}
\usepackage{float}
\usepackage{setspace}
\usepackage{parskip}
\usepackage{fancyhdr}
\usepackage{textcomp}
\usepackage{xcolor}
\usepackage{caption}
\usepackage{fvextra}
\usepackage[hidelinks]{hyperref}

\onehalfspacing
\setlength{\parskip}{6pt}
\sloppy
\setlength{\emergencystretch}{3em}
\setlength{\headheight}{40pt}
\captionsetup{font=small,labelfont=bf}
\setlength{\tabcolsep}{6pt}

\pagestyle{fancy}
\fancyhf{}
\fancyhead[L]{\includegraphics[width=1cm]{/home/yovick/repos/UAZ-IEI/Assets/UIE.png}}
\fancyhead[R]{\small Tecnologías Celulares}
\fancyfoot[C]{\thepage}
\renewcommand{\headrulewidth}{0.4pt}

\DefineVerbatimEnvironment{codigo}{Verbatim}{%
  breaklines=true,breakanywhere=true,fontsize=\small,frame=single,
  framesep=3pt,rulecolor=\color{gray!50}}

% Figura centrada con pie: \figura[ancho]{archivo}{leyenda}
\newcommand{\figura}[3][0.85]{%
  \begin{center}
    \includegraphics[width=#1\linewidth]{#2}\\[2pt]
    {\small #3}
  \end{center}}

% Portada: \portada{numero}{titulo de la practica}
\newcommand{\portada}[2]{%
\begin{titlepage}
  \thispagestyle{empty}
  \begin{center}
    \vspace*{1cm}
    \includegraphics[width=4.5cm]{/home/yovick/repos/UAZ-IEI/Assets/uaz.png}
    \vspace{0.8cm}

    {\large Universidad Autónoma de Zacatecas}\\[4pt]
    {\large Unidad Académica de Ingeniería Eléctrica}\\[4pt]
    {\normalsize Ingeniería Electrónica Industrial con especialidad en Telecomunicaciones}

    \vspace{1.8cm}
    {\LARGE\bfseries Práctica #1}\\[10pt]
    {\Large\bfseries #2}
  \end{center}

  \vspace{2.2cm}
  \begin{center}
    \begin{tabular}{rl}
      \textbf{Materia:} & Tecnologías Celulares\\
      \textbf{Docente:} & Dr. Víktor Iván Rodríguez Abdalá\\
      \textbf{Alumno:}  & Arnold Yovick Ríos Zaldívar\\
      \textbf{Fecha:}   & Noviembre de 2026\\
    \end{tabular}
  \end{center}
  \vfill
\end{titlepage}}

\usepackage[utf8]{inputenc}
\usepackage[T1]{fontenc}
\usepackage[spanish,mexico]{babel}
\usepackage{graphicx}
\usepackage[margin=2.5cm]{geometry}
\usepackage{amsmath}
\usepackage{amssymb}
\usepackage{booktabs}
\usepackage{tabularx}
\usepackage{array}
\usepackage{float}
\usepackage{setspace}
\usepackage{parskip}
\usepackage{fancyhdr}
\usepackage{textcomp}
\usepackage{xcolor}
\usepackage{caption}
\usepackage{fvextra}
\usepackage[hidelinks]{hyperref}

\onehalfspacing
\setlength{\parskip}{6pt}
\sloppy
\setlength{\emergencystretch}{3em}
\setlength{\headheight}{40pt}
\captionsetup{font=small,labelfont=bf}
\setlength{\tabcolsep}{6pt}

\pagestyle{fancy}
\fancyhf{}
\fancyhead[L]{\includegraphics[width=1cm]{/home/yovick/repos/UAZ-IEI/Assets/UIE.png}}
\fancyhead[R]{\small Tecnologías Celulares}
\fancyfoot[C]{\thepage}
\renewcommand{\headrulewidth}{0.4pt}

\DefineVerbatimEnvironment{codigo}{Verbatim}{%
  breaklines=true,breakanywhere=true,fontsize=\small,frame=single,
  framesep=3pt,rulecolor=\color{gray!50}}

% Figura centrada con pie: \figura[ancho]{archivo}{leyenda}
\newcommand{\figura}[3][0.85]{%
  \begin{center}
    \includegraphics[width=#1\linewidth]{#2}\\[2pt]
    {\small #3}
  \end{center}}

% Portada: \portada{numero}{titulo de la practica}
\newcommand{\portada}[2]{%
\begin{titlepage}
  \thispagestyle{empty}
  \begin{center}
    \vspace*{1cm}
    \includegraphics[width=4.5cm]{/home/yovick/repos/UAZ-IEI/Assets/uaz.png}
    \vspace{0.8cm}

    {\large Universidad Autónoma de Zacatecas}\\[4pt]
    {\large Unidad Académica de Ingeniería Eléctrica}\\[4pt]
    {\normalsize Ingeniería Electrónica Industrial con especialidad en Telecomunicaciones}

    \vspace{1.8cm}
    {\LARGE\bfseries Práctica #1}\\[10pt]
    {\Large\bfseries #2}
  \end{center}

  \vspace{2.2cm}
  \begin{center}
    \begin{tabular}{rl}
      \textbf{Materia:} & Tecnologías Celulares\\
      \textbf{Docente:} & Dr. Víktor Iván Rodríguez Abdalá\\
      \textbf{Alumno:}  & Arnold Yovick Ríos Zaldívar\\
      \textbf{Fecha:}   & Noviembre de 2026\\
    \end{tabular}
  \end{center}
  \vfill
\end{titlepage}}


\begin{document}

\portada{10}{Detección de símbolos en sistemas MIMO utilizando ZF y MMSE}

\textbf{Script:} \texttt{practica10.py} (Python con NumPy)

\section{Objetivo}

Implementar los detectores Zero-Forcing (ZF) y Minimum Mean Square Error (MMSE) para un
sistema MIMO 2x2 con modulación BPSK y comparar su desempeño mediante simulaciones de la
tasa de error de bit.

\section{Fundamento teórico}

El modelo del sistema es
\begin{equation}
  \mathbf{y} = \mathbf{H}\mathbf{x} + \mathbf{n},
  \qquad
  \mathbf{H}\in\mathbb{C}^{N_r\times N_t}.
\end{equation}
Cada antena transmisora envía un símbolo BPSK $\pm1$ y el receptor observa una mezcla de
ambos flujos más ruido. Los detectores lineales estiman $\mathbf{x}$ con una matriz
$\mathbf{W}$:
\begin{equation}
  \hat{\mathbf{x}} = \mathbf{W}\mathbf{y}.
\end{equation}

\paragraph{Zero-Forcing.} Anula por completo la interferencia entre flujos:
\begin{equation}
  \mathbf{W}_{\mathrm{ZF}} = \left(\mathbf{H}^{H}\mathbf{H}\right)^{-1}\mathbf{H}^{H}.
\end{equation}
Su desventaja es que amplifica el ruido cuando $\mathbf{H}$ está mal condicionada: la
potencia de ruido a la salida es proporcional a
$\left(\mathbf{H}^{H}\mathbf{H}\right)^{-1}$.

\paragraph{MMSE.} Minimiza el error cuadrático medio y balancea interferencia y ruido:
\begin{equation}
  \mathbf{W}_{\mathrm{MMSE}} = \left(\mathbf{H}^{H}\mathbf{H}
      + \frac{1}{\rho}\mathbf{I}\right)^{-1}\mathbf{H}^{H},
\end{equation}
con $\rho$ la SNR lineal. A alta SNR el término $1/\rho$ se vuelve despreciable y MMSE
tiende a ZF; a baja SNR prefiere no invertir los modos débiles del canal.

\section{Desarrollo experimental}

\subsection{Parámetros}

\begin{tabularx}{\linewidth}{@{}*{2}{>{\raggedright\arraybackslash}X}@{}}
  \toprule
  Parámetro & Valor \\
  \midrule
  Antenas & $N_t=N_r=2$ \\
  Modulación & BPSK ($\pm1$) \\
  Canal & Rayleigh i.i.d., $\mathcal{CN}(0,1)$ \\
  Bits por realización & 10\,000 por antena \\
  Realizaciones por punto & 500 \\
  $E_b/N_0$ & 0 a 20 dB en pasos de 2 dB \\
  Varianza del ruido & $\sigma^2=1/\rho$ (potencia unitaria por antena) \\
  \bottomrule
\end{tabularx}

\subsection{Procedimiento}

\begin{enumerate}
  \item Se generan bits independientes para las dos antenas y se convierten a $\pm1$.
  \item Se genera una matriz $\mathbf{H}$ Rayleigh por realización y se calcula
        $\mathbf{y}=\mathbf{H}\mathbf{x}+\mathbf{n}$.
  \item Se aplican los detectores ZF y MMSE a la misma observación.
  \item Se decide por el signo de la parte real y se acumulan los errores de ambos
        detectores.
\end{enumerate}

\section{Código}

\begin{codigo}
w_zf = np.linalg.inv(h.conj().T @ h) @ h.conj().T
w_mmse = np.linalg.inv(h.conj().T @ h + (1 / rho) * np.eye(NT)) @ h.conj().T

err_zf += np.sum((np.real(w_zf @ y) > 0).astype(int) != bits)
err_mmse += np.sum((np.real(w_mmse @ y) > 0).astype(int) != bits)
\end{codigo}

\section{Resultados}

\subsection{BER versus $E_b/N_0$}

\begin{tabular}{@{}ccc@{}}
  \toprule
  $E_b/N_0$ (dB) & BER ZF & BER MMSE \\
  \midrule
  0  & 0.15521 & 0.09651 \\
  2  & 0.11223 & 0.06692 \\
  4  & 0.07886 & 0.04462 \\
  6  & 0.04251 & 0.02476 \\
  8  & 0.03649 & 0.01838 \\
  10 & 0.02507 & 0.01221 \\
  12 & 0.01741 & 0.00827 \\
  14 & 0.01017 & 0.00518 \\
  16 & 0.00837 & 0.00428 \\
  18 & 0.00221 & 0.00166 \\
  20 & 0.00148 & 0.00102 \\
  \bottomrule
\end{tabular}

\figura{figs/ber.png}{Curvas de BER de los detectores ZF y MMSE.}

MMSE es mejor en todo el rango simulado, con una ventaja máxima a baja SNR: a 0 dB reduce
la BER de 0.155 a 0.097 y a 10 dB de 0.0251 a 0.0122 (factor 2.05). A 18--20 dB las curvas
se acercan, como predice la teoría, porque el término de ruido de la matriz MMSE se vuelve
despreciable.

\subsection{Constelaciones recuperadas}

\figura{figs/constelaciones.png}{Símbolos detectados por ZF y MMSE a $E_b/N_0=10$ dB (primer flujo).}

Con ZF la nube es más ancha y aparecen puntos alejados del eje real: la inversión del canal
amplificó el ruido en las direcciones débiles. Con MMSE la dispersión es menor, sobre todo
en los símbolos que caen cerca del origen.

\section{Análisis de resultados}

\begin{itemize}
  \item \textbf{¿Por qué MMSE suele superar a ZF?} Porque no invierte por completo el canal:
        añade un término de regularización $1/\rho$ que evita amplificar el ruido en los
        modos con poca ganancia. A cambio acepta algo de interferencia residual.
  \item \textbf{¿En qué condiciones ambos detectores presentan desempeños similares?} A alta
        SNR, cuando $1/\rho$ es pequeño frente a los valores propios de
        $\mathbf{H}^{H}\mathbf{H}$; en ese límite MMSE converge a ZF.
  \item \textbf{¿Cómo afecta el ruido a cada detector?} ZF elimina la interferencia pero
        amplifica el ruido según el condicionamiento de $\mathbf{H}$; MMSE reparte el error
        entre interferencia y ruido, por lo que es más robusto cuando el canal está mal
        condicionado.
  \item \textbf{¿Qué ventajas ofrece cada método?} ZF es simple y elimina totalmente la
        interferencia entre flujos; MMSE es óptimo en el sentido de mínimo error cuadrático
        medio y funciona mejor a baja SNR. Ambos son lineales y de bajo costo comparados
        con detectores no lineales como ML o V-BLAST.
\end{itemize}

\section{Cuestionario}

\begin{enumerate}
  \item \textbf{¿Qué hace el detector ZF?} Aplica la pseudo-inversa del canal para anular la
        interferencia entre los flujos transmitidos y recuperar cada símbolo por separado.
  \item \textbf{¿Qué desventaja presenta ZF?} Amplifica el ruido cuando la matriz de canal
        está mal condicionada o tiene desvanecimientos profundos, lo que degrada la BER a
        baja SNR.
  \item \textbf{¿Qué ventaja tiene MMSE?} Balancea la interferencia y el ruido mediante la
        regularización $1/\rho$, por lo que ofrece menor BER que ZF, especialmente a baja
        SNR.
  \item \textbf{¿Por qué MMSE considera la potencia de ruido?} Porque su criterio es
        minimizar $\mathbb{E}\{|\hat{\mathbf{x}}-\mathbf{x}|^2\}$, que depende de la
        varianza del ruido; el término $1/\rho$ es precisamente la inversa de la SNR.
  \item \textbf{¿Qué detector utiliza LTE con técnicas avanzadas de recepción?} LTE usa
        ecualización MMSE (y variantes como MMSE-IRC para interferencia) en el receptor,
        junto con cancelación sucesiva de interferencia y detección ML en esquemas
        avanzados.
\end{enumerate}

\section{Conclusiones}

Se implementaron los detectores ZF y MMSE para un MIMO 2x2 con BPSK sobre canal Rayleigh.
MMSE presentó menor BER en todo el rango, con una mejora de 2.05 veces a 10 dB
(0.0122 frente a 0.0251) y una diferencia que se reduce a alta SNR. Las constelaciones
mostraron la amplificación de ruido característica de ZF. El experimento confirma que la
regularización de MMSE es la opción práctica cuando la SNR no es alta.

\end{document}
